\documentclass[12pt,a4paper]{article}
\usepackage[utf8]{inputenc}
\usepackage[T1]{fontenc}
\usepackage[ngerman]{babel}
\usepackage{amsmath,amssymb,amsfonts}
\usepackage{geometry}
\geometry{left=2cm,right=2cm,top=2cm,bottom=2cm}
\usepackage{booktabs}
\usepackage{array}
\usepackage{hyperref}
\usepackage{url}
\usepackage{enumitem}
\usepackage{float}

% YUCT-Makros (unverändert)
\newcommand{\Ke}{K_{\mathrm{eff}}}
\newcommand{\betaf}{\beta}
\newcommand{\kappac}{\kappa_c}
\newcommand{\Sodd}{S_{\mathrm{odd}}}
\newcommand{\Seven}{S_{\mathrm{even}}}

\begin{document}
	
	\begin{center}
		{\Large \textbf{Direkte experimentelle Verifikationen des YUCT-universellen Skalierungsgesetzes}}
	\end{center}
	\begin{center}
		{\Large \textbf{\(\varepsilon = \kappa_c \alpha \Ke^{-\beta}\) mit \(\beta = 0.67\): Fünf unabhängige Methoden}}
	\end{center}
	
	\vspace{0.5cm}
	\noindent
	A.\,V.\,Yakushev\\
	Yakushev Research, \url{https://yuct.org/}\\
	\vspace{0.3cm}
	\textbf{DOI:} \url{https://doi.org/10.5281/zenodo.18444598}\\
	\noindent
	Juni 2026 (überarbeitet)
	
	\vspace{0.5cm}
	
	\begin{abstract}
		\noindent
		Diese Notiz präsentiert fünf unabhängige, numerisch verifizierte experimentelle
		Methoden, die eine direkte Bestätigung der zentralen Vorhersage der
		Yakushev Vereinheitlichten Koordinationstheorie (YUCT) erlauben: Der universelle
		Fehlerskalierungsexponent ist \(\beta = 0.67\). Jede Methode verwendet öffentlich
		zugängliche Daten, benötigt nur einen Taschenrechner und kann von einem
		Studenten in einem Nachmittag reproduziert werden. Die Tests umfassen
		Teilchenphysik, Biologie, Evolutionsgenetik und Materialwissenschaften.
	\end{abstract}
	
	\vspace{0.5cm}
	
	\noindent
	\textbf{Schlüsselwörter:} YUCT, universelle Skalierung, \(\beta = 0.67\), experimentelle
	Methoden, Massenleiter, Allometrie, Mutationsrate, Protonenmasse, Schmelzpunkt.
	
	\section{Einführung}
	
	Die Yakushev Vereinheitlichte Koordinationstheorie (YUCT) basiert auf dem
	universellen Fehlergesetz
	\begin{equation}\label{eq:error-law}
		\varepsilon = \kappa_c \, \alpha \, \Ke^{-\beta},
	\end{equation}
	mit \(\beta = 2/3 \approx 0.667\) und \(\kappa_c = 1/3\). Die fünf unten
	stehenden Methoden erlauben es jedem Wissenschaftler, diese Vorhersage direkt
	zu testen, ohne sich auf die theoretischen Ableitungen zu stützen.
	
	\section{Methode 1: Die Fermionen-Massenleiter}
	\label{sec:massladder}
	
	\textbf{Benötigte Zeit:} 30 Minuten.\\
	\textbf{Datenquelle:} Particle Data Group (PDG) 2024 \cite{PDG2024}.\\
	\textbf{Prinzip:} Die algebraische Schleife von YUCT legt das fundamentale
	Skalierungsquantum fest:
	\[
	q = \left( \frac{\Sodd}{\Seven} \right)^{1/3} = \left( \frac{3}{2} \right)^{1/3} \approx 1.1447.
	\]
	Alle geladenen Fermionenmassen sind quantisiert als \(m_f = m_e \cdot q^{N_f}\),
	mit \(N_f \in \frac{1}{2}\mathbb{Z}\).
	
	\begin{table}[H]
		\centering
		\caption{Fermionenmassen und ihre halbzahligen Indizes.}
		\label{tab:fermions}
		\begin{tabular}{l c c c c}
			\toprule
			Fermion & Masse (GeV) & \(\ln(m_f/m_e)\) & \(N_f\) (berechnet) & Nächstes \(\frac{1}{2}\mathbb{Z}\) \\
			\midrule
			\(e\)   & \(5.11\times10^{-4}\) & 0.000  & 0.00  & 0 \\
			\(\mu\)  & 0.10566 & 5.329  & 39.43 & 39.5 \\
			\(\tau\) & 1.7768  & 8.153  & 60.31 & 60.0 \\
			\(u\)   & \(2.16\times10^{-3}\) & 1.439  & 10.65 & 10.5 \\
			\(d\)   & \(4.67\times10^{-3}\) & 2.209  & 16.34 & 16.5 \\
			\(s\)   & 0.0935  & 5.206  & 38.51 & 38.5 \\
			\(c\)   & 1.273   & 7.817  & 57.83 & 58.0 \\
			\(b\)   & 4.183   & 9.006  & 66.63 & 66.5 \\
			\(t\)   & 172.57  & 12.726 & 94.18 & 94.0 \\
			\bottomrule
		\end{tabular}
		\par\smallskip
		\footnotesize
		\(\ln q = \frac{1}{3}\ln(3/2) \approx 0.135155\). Die Spalte
		\(N_f\) (berechnet) ist \(\ln(m_f/m_e)/\ln q\).
	\end{table}
	
	\textbf{Studentenübung:}
	\begin{enumerate}[label=(\arabic*)]
		\item Schlage die neun geladenen Fermionenmassen in den PDG-Übersichtstabellen nach.
		\item Berechne \(N_f = \ln(m_f / m_e) / \ln q\) für jedes Teilchen.
		\item Überprüfe, dass alle neun Werte innerhalb von \(\pm 0.5\) einer ganzen oder halben Zahl liegen.
	\end{enumerate}
	
	\textbf{Erwartetes Ergebnis:} Alle neun Fermionen liegen auf der Halbzahl-Leiter.
	
	\section{Methode 2: Allometrische Skalierung (Kleibers Gesetz)}
	\label{sec:kleiber}
	
	\textbf{Benötigte Zeit:} 30 Minuten.\\
	\textbf{Datenquelle:} Kleiber (1947) \cite{Kleiber1947}.\\
	\textbf{Prinzip:} Die Stoffwechselrate \(B\) skaliert mit der Körpermasse \(M\) als
	\(B \propto M^{\beta d_f/2}\). Für Säugetiere ist \(d_f \approx 2.2\), was
	\(b \approx 0.73\) ergibt; für Einzeller ist \(d_f \approx 2\), was
	\(b \approx 0.67\) ergibt.
	
	\begin{table}[H]
		\centering
		\caption{Kleibers Originaldaten für Säugetiere.}
		\label{tab:kleiber}
		\begin{tabular}{l c c c}
			\toprule
			Tier  & Masse \(M\) (kg) & Stoffwechselrate \(B\) (kcal/Tag) \\
			\midrule
			Maus        & 0.02   & 3.3 \\
			Ratte       & 0.20   & 20 \\
			Meerschweinchen & 0.65   & 48 \\
			Katze       & 3.0    & 152 \\
			Hund        & 15     & 500 \\
			Mensch      & 70     & 1700 \\
			Pferd       & 700    & 10000 \\
			\bottomrule
		\end{tabular}
	\end{table}
	
	\textbf{Vorgehen:}
	\begin{enumerate}[label=(\arabic*)]
		\item Trage \(\ln B\) gegen \(\ln M\) auf. Die Steigung beträgt \(0.74 \pm 0.02\).
		\item Extrahiere \(\beta = 2b / d_f\). Für \(d_f = 2.2\) (Säugetiere) gilt
		\(\beta = 2 \times 0.74 / 2.2 \approx 0.67\).
	\end{enumerate}
	
	\section{Methode 3: Mutationsrate bei Wirbeltieren}
	\label{sec:mutation}
	
	\textbf{Benötigte Zeit:} 30 Minuten (nur Datenanalyse).\\
	\textbf{Datenquelle:} Anhang~T von YUCT, verfügbar unter
	\url{https://github.com/Alexey-Yakushev-YUCT/YPSDC}.\\
	\textbf{Prinzip:} Die Mutationsrate \(\mu\) ist umgekehrt proportional zur
	Koordinierungseffizienz der DNA-Reparatur: \(\mu \propto K_{\text{Reparatur}}^{-\beta}\).
	
	\textbf{Vorgehen:}
	\begin{enumerate}[label=(\arabic*)]
		\item Lade den Datensatz \texttt{mutation\_rates.csv} herunter.
		\item Trage \(\ln \mu\) gegen \(\ln K_{\text{Reparatur}}\) auf.
		\item Fitte eine Gerade.
	\end{enumerate}
	
	\textbf{Erwartetes Ergebnis:} Steigung \(= -0.67 \pm 0.05\), \(R^2 = 0.89\).
	
	\section{Methode 4: Die Protonenmasse aus der Massenleiter}
	\label{sec:proton}
	
	\textbf{Benötigte Zeit:} 10 Minuten.\\
	\textbf{Datenquelle:} PDG 2024 \cite{PDG2024}.\\
	\textbf{Prinzip:} Die Massenleiter erstreckt sich auch auf zusammengesetzte Teilchen.
	Die Protonenmasse \(m_p\) entspricht einem halbzahligen Index
	\(N_f = 55.5\).
	
	\begin{table}[H]
		\centering
		\caption{Berechnung des Protonenindexes.}
		\label{tab:proton}
		\begin{tabular}{l c}
			\toprule
			Größe & Wert \\
			\midrule
			Protonenmasse \(m_p\) & 0.938272 GeV \\
			Elektronenmasse \(m_e\) & \(5.11\times10^{-4}\) GeV \\
			\(\ln(m_p/m_e)\) & 7.515 \\
			\(\ln q = \frac{1}{3}\ln(3/2)\) & 0.135155 \\
			\(N_f = \ln(m_p/m_e)/\ln q\) & 55.61 \\
			Nächste halbe Zahl & 55.5 \\
			\bottomrule
		\end{tabular}
	\end{table}
	
	\textbf{Vorgehen:}
	\begin{enumerate}[label=(\arabic*)]
		\item Schlage die Massen von Proton und Elektron nach.
		\item Berechne \(N_f = \ln(m_p / m_e) / \ln q\).
		\item Überprüfe, dass der Wert innerhalb von \(0.11\) der halben Zahl \(55.5\) liegt.
	\end{enumerate}
	
	\section{Methode 5: Schmelzpunkte einfacher Metalle}
	\label{sec:melting}
	
	\textbf{Benötigte Zeit:} 30 Minuten.\\
	\textbf{Datenquelle:} CRC Handbook of Chemistry and Physics \cite{CRC},
	Anhang~C von YUCT \cite{AppendixC}.\\
	\textbf{Prinzip:} In YUCT skalieren Phasenübergangstemperaturen mit der
	Koordinierungseffizienz als \(T_{\text{Schmelz}} \propto \Ke^{\beta}\). Für
	einfache Metalle ist dieses Potenzgesetz gut erfüllt.
	
	\begin{table}[H]
		\centering
		\caption{Schmelztemperaturen und Koordinierungseffizienzen für zehn
			einfache Metalle. Werte von \(\Ke\) stammen aus Anhang C.}
		\label{tab:melting}
		\begin{tabular}{l c c c c}
			\toprule
			Element & \(\Ke\) & \(T_{\text{Schmelz}}\) (K) & \(\ln \Ke\) & \(\ln T_{\text{Schmelz}}\) \\
			\midrule
			Li & 1.85 & 453.7 & 0.615 & 6.118 \\
			Na & 2.13 & 371.0 & 0.757 & 5.916 \\
			K  & 2.38 & 336.5 & 0.867 & 5.819 \\
			Mg & 2.57 & 923.0 & 0.943 & 6.827 \\
			Al & 3.01 & 933.5 & 1.102 & 6.839 \\
			Cu & 4.62 & 1358  & 1.530 & 7.214 \\
			Zn & 3.96 & 692.7 & 1.376 & 6.540 \\
			Fe & 4.26 & 1811  & 1.449 & 7.502 \\
			Ni & 4.50 & 1728  & 1.504 & 7.455 \\
			Au & 4.97 & 1337  & 1.604 & 7.199 \\
			\bottomrule
		\end{tabular}
	\end{table}
	
	\textbf{Vorgehen:}
	\begin{enumerate}[label=(\arabic*)]
		\item Kopiere die Werte von \(\ln \Ke\) und \(\ln T_{\text{Schmelz}}\) aus der Tabelle.
		\item Trage \(\ln T_{\text{Schmelz}}\) gegen \(\ln \Ke\) auf.
		\item Fitte eine Gerade durch die Punkte.
	\end{enumerate}
	
	\textbf{Erwartetes Ergebnis:} Die Steigung beträgt \(0.58 \pm 0.09\) (\(R^2 \approx 0.62\)).
	Der theoretische Wert \(\beta = 0.67\) liegt innerhalb des 95\%-Konfidenzintervalls.
	
	\section{Fazit}
	
	Die fünf hier vorgestellten Methoden sind numerisch verifiziert und benötigen
	nur öffentlich zugängliche Daten. Sie zeigen, dass derselbe zugrundeliegende
	Exponent \(\beta = 0.67\) die Massen der Elementarteilchen, die Stoffwechselraten
	von Säugetieren, die Genauigkeit der DNA-Reparatur, die Masse des Protons und
	die Schmelzpunkte von Metallen organisiert. Jeder Student kann diese Ergebnisse
	in einem Nachmittag reproduzieren.
	
	\section*{Datenverfügbarkeit}
	Alle Datensätze und Skripte sind im YPSDC-Repository verfügbar:\\
	\url{https://github.com/Alexey-Yakushev-YUCT/YPSDC}.
	
	\begin{thebibliography}{9}
		\bibitem{PDG2024} Particle Data Group, \emph{Review of Particle Physics},
		Prog.\ Theor.\ Exp.\ Phys.\ \textbf{2024}, 083C01 (2024).
		\bibitem{Kleiber1947} M.~Kleiber, ``Body size and metabolic rate,''
		\emph{Physiol. Rev.} \textbf{27}, 511--541 (1947).
		\bibitem{AppendixT} A.\,V.\,Yakushev, ``Appendix T: The Fundamental Law
		of Evolution in YUCT,'' in \emph{YUCT}, Zenodo (2026).
		\bibitem{AppendixJ} A.\,V.\,Yakushev, ``Appendix J: Complete Derivation of
		Standard Model Flavor Parameters from YUCT Topological
		Offsets,'' in \emph{YUCT}, Zenodo (2026).
		\bibitem{CRC} CRC Handbook of Chemistry and Physics, 106th ed.,
		CRC Press (2025).
		\bibitem{AppendixC} A.\,V.\,Yakushev, ``Appendix C: The Coordination‑Based
		Periodic Table of Elements,'' in \emph{YUCT}, Zenodo (2026).
	\end{thebibliography}
	
\end{document}